\chapter{Pregunta de Investigación}

¿En qué medida la integración de un modelo matemático basado en ecuaciones diferenciales ordinarias, sensores IoT de bajo costo y algoritmos de aprendizaje automático permite desarrollar un sistema híbrido que no solo estime con precisión las emisiones de \ce{CO2} y \ce{CH4}, sino que constituya una herramienta de monitoreo dinámico para fundamentar la toma de decisiones operativas y generar métricas verificables del desempeño ambiental en la bioconversión de residuos agroindustriales con \textit{Hermetia illucens}?

\chapter{Hipótesis de Investigación}

\section{Hipótesis General}

Un sistema híbrido que integre modelamiento matemático basado en ecuaciones diferenciales ordinarias (EDO), sensores IoT de bajo costo y algoritmos de aprendizaje automático permite estimar las emisiones de \ce{CO2} y \ce{CH4} en tiempo real con un error absoluto medio porcentual (MAPE) inferior al 10\,\% respecto a mediciones de referencia. Esta capacidad de monitoreo continuo proporciona, a su vez, indicadores dinámicos sobre la eficiencia metabólica del proceso, permitiendo identificar ineficiencias operativas (como la anaerobiosis) y cuantificar el desempeño ambiental del sistema más allá de las estimaciones estáticas tradicionales.

\section{Hipótesis Específicas}
\subsection{H1: Precisión instrumental de bajo costo}
Los sensores NDIR de bajo costo, sometidos a protocolos de calibración y preacondicionamiento de señal, permiten medir concentraciones de \ce{CO2} y \ce{CH4} en el entorno de bioconversión con una desviación (MAPE) $\leq 10$\,\% respecto a equipos certificados de laboratorio.

\subsection{H2: Capacidad predictiva del modelo mecanístico}
El modelo de EDO, fundamentado en la bioenergética del crecimiento larval y parametrizado con datos locales, es capaz de describir la tendencia base de las emisiones de \ce{CO2} y \ce{CH4} en ventanas temporales de 30 minutos con un MAPE $\leq 15$\,\%, capturando la dinámica biológica subyacente.

\subsection{H3: Refinamiento por aprendizaje automático}
La integración de una capa de corrección mediante algoritmos de aprendizaje automático reduce el error residual del modelo de EDO en al menos un 30\,\%, logrando que el sistema híbrido alcance un MAPE $< 10$\,\% ante perturbaciones no modeladas.
\subsection{H4: Valor para la gestión operativa}
El sistema híbrido integrado es capaz de identificar patrones de emisión anómalos (como picos imprevistos de \ce{CH4} indicativos de anaerobiosis) que no son detectables mediante modelado estático, proporcionando indicadores de desempeño ambiental en tiempo real útiles para la toma de decisiones de manejo.

\emph{Nota: Las métricas de precisión (H1--H3) se evaluarán bajo condiciones experimentales controladas en laboratorio (30\,°C $\pm$ 2\,°C, HR 60--70\,\%), mientras que la H4 se validará mediante análisis de escenarios operativos.}